% page 392 du fac-simile (image p-391 du scan)
% bloc 157.5 x 237.7 mm ; marge gauche 8.0 mm
\pgc{-12.700mm}{0.000mm}{
\l{\cel{3}384}
\l{}
\l{moruo. (\sou{zool.}) Marfisho ek la genero \textquotedbl{}gadi\textquotedbl{}, de qui un}
\l{\cel{3}speco, qua tre abundas en Nov-Lando, esas la objekto}
\l{\cel{3}di komerco tre granda. - L.\cel{1}\sou{gadus morrhus}. - FL.}
\l{}
\l{moruso. (\sou{bot.)}\cel{2}Frukto monoika, pulpoza, di arboro kun}
\l{\cel{3}folii alternanta. - L.\cel{1}\sou{morus}.\cel{2}- SL.}
\l{}
\l{\sou{mor}vo. (\sou{patol.}) Morbo kontagiala di la kavali, maxim-}
\l{\cel{3}multa-kaze mortiganta, karakterizata da la inflameso}
\l{\cel{3}e la ekfluado di la membrano pituitala. - FS.}
\l{}
\l{\sou{moso}. (\sou{navig.}) Adolecanto qua, sur navo, lernas la mestiero}
\l{\cel{3}di maristo (+kruano). - FI.}
\l{}
\l{\sou{moskardino}. (\sou{zool.)}\cel{2}Speco di gliro, di qua la odoro for-}
\l{\cel{3}ta memorigas musko. - L.\cel{1}\sou{muscardinus avelianarius}.}
\l{}
\l{\sou{mosko}. Substanco odoranta quan kontenas posho renkontrebla}
\l{\cel{3}sur la ventro di la maskulo di sorto di kapreol-yuno.DEFIR.}
\l{}
\l{\sou{moskeo.(religio)}. Templo di la Mohamedisti. - DEFIRS.}
\l{}
\l{\sou{moskito}. (\sou{zool.)}\cel{2}Insekto diptera, di qua la pikuro esas do-}
\l{\cel{3}loriganta. - DEFRS.}
\l{}
\l{\sou{mosto}.\cel{2}Suko di vitbero, pronta subisar la fermentaco}
\l{\cel{3}alkoholifera. - DEFIS.}
\l{}
\l{\sou{moto}. Sentenco, proverbo preferata da ulu, od atribuita ad}
\l{\cel{3}ulu, kom expresanta plu o min adequate lua maniero pen-}
\l{\cel{3}sar, sentar. - DEIR.}
\l{}
\l{\sou{motacilo}. (\sou{zool.}) Mikra ucelo qua flugetas cirkum la bestio-}
\l{\cel{3}trupi, ed apud la aqui fluanta o dormanta. - L.\cel{1}\sou{mola}-}
\l{\cel{3}\sou{cilla alba.}\cel{2}L.}
\l{}
\l{\sou{moteto.}\cel{2}I. (\sou{olim}). Kirko-kantajo, mi-sakra e mi-profana,}
\l{\cel{3}plura-parta, tolerita apud la plan-kanto. - II. Muzikajo}
\l{\cel{3}destinita a la kirko, e kompozita ek vorti Latina quin}
\l{\cel{3}onu prenis exter la oficio, quale psalmo, himno, anti-}
\l{\cel{3}fono, e c. DEFIRS.}
\l{}
\l{\sou{motivo}. I. La \textquotedbl{}pro quo\textquotedbl{} o \textquotedbl{}por quo\textquotedbl{} di ago o di faco. -}
\l{\cel{3}II. En arto-verko : la temo developenda. - DEFIRS.}
\l{}
\l{\sou{motoro}. Omna sistemo materiala destinita a transformar}
\l{\cel{3}irga energio aden energio mekanikala. - DEFIRS.}
\l{}
\l{\sou{movar}. (\sou{trans}.)I. Igar funcionar. (\sou{anke metaf}.). II.}
\l{\cel{3}(pri korpo o parto di korpo). Diplasar su. EFIS.}
\l{}
\l{\sou{movimento}. (\sou{muziko)}. Rapideso-grado ye qua frazo muzikala}
\l{\cel{3}pleesas o kantesas. - EFI.}
}
